
Weight-space encoders that predict properties of trained neural networks from their parameters face a structural challenge: fully-connected networks are invariant to neuron permutations within layers, yet standard flat-MLP encoders treat neuron position as input signal. This work presents a controlled comparison between Neural Functional Transformers (NFT), a permutation-equivariant encoder, and flat-MLP baselines for generalization gap prediction on the Unterthiner model zoo (29,997 models). In permutation stress tests at severity levels s in {0, 0.25, 0.5, 1.0}, flat-MLP exhibited a Spearman rank correlation drop of 0.160 (from 0.303 to 0.143; bootstrap p = 0.000), while NFT showed no measurable degradation (delta-rho = 4.09 x $10^-6$; bootstrap p = 0.477). A six-encoder ablation (18 training runs, 3 seeds) confirmed larger flat-MLP degradation (delta-rho = 0.640) and near-zero NFT sensitivity (delta-rho = 4.71 x $10^-7$). Mediation analysis yielded delta-R-squared = 0.228, indicating that NFT embeddings explained additional variance in generalization gap prediction beyond augmentation-based approaches. Permutation augmentation reduced flat-MLP sensitivity by approximately 65\% on average but exhibited high seed-dependent variance (delta-rho range: 0.096--0.317). L2-norm canonicalization collapsed to constant predictions across all seeds. NFT used 75K parameters compared to flat-MLP's 3.04M. Two of five planned sub-hypotheses (cross-pipeline transfer, graceful degradation curves) were not executed; cross-pipeline claims remain untested.
